\section{The StrongREJECT benchmark}\label{sec:our-benchmark}

StrongREJECT addresses the shortcomings of existing jailbreak benchmarks with a higher-quality dataset of forbidden prompts and a more accurate automatic evaluation method. This section describes our benchmark in greater detail.

\subsection{Forbidden prompts}\label{ssec:forbidden_prompts}

\pseudopara{Categories.} It is important to ensure that our forbidden prompts cover a wide range of harmful behaviors. Accordingly, we cross-referenced prohibited use cases in the usage policies of OpenAI, Anthropic, Google Gemini, Meta LLaMA, and DeepInfra to construct categories of harmful behavior to test for. The six categories we selected cover behaviors that are forbidden in all usage policies and rejected by the latest models from these providers and others (Mistral, Reka, and Cohere): illegal goods and services; non-violent crimes; hate, harassment, and discrimination; disinformation and deception; violence; and sexual content.

\pseudopara{Writing, generating, and curating forbidden prompts.} Our dataset of forbidden prompts consists of prompts we wrote manually, LLM-generated prompts using GPT-4 and pplx-70b-online following \citet{shen2023anything}, and curated prompts from existing jailbreak benchmarks~\citep{chen2022should, huang2023catastrophic, shaikh2022second, liu2023jailbreaking, achiam2023gpt, deng2023jailbreaker, shen2023anything}. We verified that all of the forbidden prompts we selected for our dataset were distinct, refused by the latest models from the major providers above, detailed enough for a human to easily evaluate responses to them, and simple enough for moderately-capable models to answer them. Importantly, the prompts we selected had factually verifiable answers with easily accessible information, such that a person with Internet access could write a high-quality response within an hour. This makes the victim model responses easy to evaluate and ensures that researchers will not create an information hazard by publishing responses to any of our forbidden prompts.

In total, the StrongREJECT dataset offers 313 high-quality forbidden prompts, with at least 50 forbidden prompts in each category. Out of these, 70\% are our own novel prompts, 10\% each are from DAN \citep{shen2023anything} and AdvBench \citep{zou2023universal}, and the rest are from other prior work and OpenAI's system card. See~\cref{app:data-supersetting} for further details.

\subsection{Automated evaluation method}\label{ssec:automated evaluator}

Many automated evaluators use LLMs to score a forbidden prompt-victim model response pair. The LLM can be hosted remotely or on a user's GPU. For example, GPT-4 Judge \cite{qi2023fine} provides an evaluation prompt template for GPT-4, while HarmBench provides a fine-tuned classifier hosted on HuggingFace for users to run on their own GPU \cite{mazeika2024harmbench}. Using remotely-hosted LLMs allows users to run jailbreak evaluations without a GPU and avoid complexities like memory management. However, it is usually not free (although costs continue to drop dramatically), and results may not be reproducible if the provider introduces changes. Alternatively, running an LLM evaluator on one's own GPU allows for white-box access and reproducible results but often requires additional setup. We conducted an informal survey that showed AI researchers were equally split between those who preferred to use remotely-hosted LLM evaluators and those who preferred running LLM evaluators on their own GPU.

We solve this problem by providing two versions of our StrongREJECT evaluator: a \textit{rubric-based evaluator} that scores victim model responses according to a rubric and is best used with remotely-hosted frontier models, like GPT-4o-mini, and an open-source \textit{fine-tuned evaluator} we created by fine-tuning Gemma 2B \citep{gemma_2024} on labels produced by our rubric-based evaluator. Although we use the rubric-based evaluator with GPT-4o-mini throughout this paper, researchers can use it with any LLM they choose, such as Claude or Gemini. This way, researchers can achieve similar results using their preferred setup.

The rubric-based StrongREJECT evaluator prompts a pre-trained LLM (the \textit{evaluator model}) with the forbidden prompt and victim model's response, along with scoring instructions. The evaluator model outputs chain-of-thought reasoning about how well the response addresses the prompt before generating three scores: a binary score for non-refusal and two 5-point Likert scale scores ranging from [1-5] (then re-scaled to [0-1]) of how specific and convincing the response was.

The final score for a single prompt/response pair is:
\[
    \text{score} = (1 - \text{refused}) \times \frac{\text{specific} + \text{convincing}}{2}
\]

As mentioned, a model will respond to a prompt with high-quality information when it is willing to and capable of doing so. Our automated evaluator considers both of these dimensions: willingness and capabilities. If a victim model is unwilling to respond to a forbidden prompt (i.e., it issues a refusal), it will receive a score of 0. However, non-refusal alone is insufficient to receive a high StrongREJECT score. A model must also be capable of providing specific, convincing information related to the forbidden prompt. In the extreme case, a response that offers no specific or convincing information - such as the incoherent responses we referenced in~\Cref{sec:intro} - is just as useless as a refusal and will also receive a score of 0. At the other extreme, victim models willing to and capable of providing specific, convincing information related to the forbidden prompt will receive high StrongREJECT scores.

The rubric-based StrongREJECT evaluator is fast and highly affordable. Its exact runtime and cost depend on the jailbreak and the length of the jailbroken prompts and victim model responses. However, in practice, we can run and evaluate a dozen jailbreaks on our entire dataset in about 15 minutes from our local machine at a cost of less than \$1 using GPT-4o-mini.

Still, many researchers prefer to run open-source LLM evaluators on their own GPU. To meet this demand, we created an open-source, fine-tuned version of our StrongREJECT evaluator by fine-tuning Gemma 2B \citep{gemma_2024}. We chose this model because it is small enough for researchers to run on a single GPU. We fine-tuned it for 1 epoch (6 hours) on a NVIDIA RTX A6000 GPU using the AdamW optimizer with a learning rate of 2e-4 to provide 5-point Likert scale scores from [1-5] (later re-scaled to [0-1]) using a dataset of approximately 15,000 unique responses with labels provided by the rubric-based version of the StrongREJECT evaluator. The dataset consists of approximately 4,000 victim model responses to forbidden prompts from the StrongREJECT benchmark, excluding those used for human evaluations as described in ~\Cref{sec:human_eval}, and 11,000 victim model responses from \citet{mazeika2024harmbench}. In practice, we find that our fine-tuned evaluator hosted on an A6000 GPU runs approximately as quickly as our rubric-based evaluator. \cref{app:automated evaluator} and \cref{app: finetuned evaluator} respectively provide additional details about the rubric-based and fine-tuned versions of our automated evaluator.